% Standalone source-record documentation; no research attempts.
\documentclass[11pt]{article}
\usepackage[margin=25mm]{geometry}
\usepackage{fontspec,unicode-math}
\setmainfont{Latin Modern Roman}
\setmathfont{Latin Modern Math}
\usepackage{amsmath,amssymb,mathtools,xurl,hyperref}
\hypersetup{colorlinks=true,urlcolor=blue,linkcolor=blue}
\setcounter{secnumdepth}{0}
\setlength{\parindent}{0pt}
\setlength{\parskip}{5pt}
\setlength{\emergencystretch}{3em}
\newcommand{\sref}[2]{\hyperlink{source:#1:#2}{[S#2]}}
\newcommand{\cataloguescope}[1]{\paragraph{Source record.}#1}
\newcommand{\floor}[1]{\left\lfloor#1\right\rfloor}
\begin{document}
\section{Exact support and trace-controlled ellipsoidal reachability}\label{statement}

\subsection{Definitions and mathematical statement}

\paragraph{Shared notation.}Standard number systems are $\mathbb N,\mathbb Z,\mathbb Q,\mathbb R,\mathbb C$. Specialized notation and index ranges follow the cited source and the definitions below.


A reachable set consists of states that can be reached by admissible controls and disturbances over a stated time horizon. Hamilton--Jacobi--Bellman equations describe optimal-control value functions; Hamilton--Jacobi--Isaacs equations add competing controls in a differential game. An ellipsoid is $\{x:(x-c)^{\mathsf T}Q^{-1}(x-c)\leq1\}$ for $Q\succ0$, and ellipsoidal calculus computes inner or outer approximations of evolving sets.
\paragraph{Record scope and qualifications.}
The catalogue title is retained as metadata. Named certificates, contracts, bounds or frameworks in that title are not asserted as results of the original source extract.
The excerpt is a broad discussion with no fixed controlled dynamics or accuracy criterion. Its mention of ellipsoidal methods does not specify exact support or trace bounds. The supplied catalogue marks this record open; its present research status was not independently checked. Its supplied reconstruction is labelled unverified.
The source text below is rendered literally where necessary; damaged formulas and composite extracts are not silently repaired into a new mathematical claim.
The supplied catalogue marks this formulation as an unverified reconstruction; missing assumptions are not certified by the source link.
\paragraph{Statement recorded in the supplied source.}
A.5 Kurzhanski, Alexander

Among the central topics of modern control theory are problems of control synthesis for complex systems. These include problems of control under uncertainty and conflict as well as those under complex constraints. The solutions to such problems are well formalizable within Hamiltonian techniques and Dynamic Programming ideas, being as a rule reduced to the solution of HJB and HJBI partial differential equations or/and variational inequalities. An example where such methods are successfully applied are problems of forward and backward reachability for uncertain systems with further passage to control synthesis, safety verification, measurement feedback and related issues. The advent of HJB methods for continuous systems is due to the introduction of appropriate theories for generalized viscosity-type solutions and their equivalents. Another important topic is the solution to problems in dynamics and control for set-valued systems in terms of Hamiltonian formalism and Dynamic Programming. At present there is a sharp necessity for a breakthrough in numerical methods for solving equations and variational inequalities of the HJB-HJBI types motivated by problems of the above. There is also an interest in numerical methods for the comparison principle in HJB theory, which allows to calculate upper and lower bounds to exact solutions of HJB equations. Another perspective is the calculation of set-valued solutions to problems in evolution dynamics, estimation and control. Together with my colleagues we have developed an ellipsoidal calculus aimed at such problems and closely connected to HJB theory. The calculus is applicable to systems with original linear structure, but its methods allow effective calculations in high dimensions as well as computer animation in high dimensions through computer windows. For nonlinear systems some comparison methods were indicated. 8

The discussion of such methods as well as those for stochastic dynamics are within my primary interests at the AIM workshop. \sref{statement}{2}
\paragraph{Source status.}
The supplied catalogue labels this source record open; that metadata was checked against the supplied snapshot on 2026-10-11. The label is preserved for provenance and is not an independent verification of present open status or mathematical correctness.
\subsection{Short English statement}
The participant discusses numerical control, reachability, comparison principles, and ellipsoidal set approximations.
\subsection{Sources}
\begin{itemize}
\item[\hypertarget{source:statement:1}{S1}] Ulam AI and the UnsolvedMath Contributors, UnsolvedMath: A Curated Collection of Open Mathematics Problems, catalogue entry. Catalogue curation and metadata: CC BY 4.0. The source extract is retained; definitions, English summaries and qualifications are editorial additions.
\url{https://huggingface.co/datasets/ulamai/UnsolvedMath}.
\item[\hypertarget{source:statement:2}{S2}] AIM Workshop Problem Lists, source linked by the catalogue.
\url{https://aimath.org/WWN/optimalcontrol/optimalcontrol.pdf}.
\end{itemize}
\paragraph{Preparation and review.}
The original source and UnsolvedMath contributors retain attribution for the source extract. Codex (GPT-6; exact serving revision not exposed) assisted with definitions, English summaries, scope notes and formatting. No new human mathematical authorship or independently verified solution is asserted. Catalogue identity checked against the complete supplied registry and current repository. Distinct upstream records, including repeated questions and broad or damaged extracts, are retained for complete coverage of the source CS category. A repeated record is not a new independent conjecture. Complete mathematical-equivalence and current-literature review has not been performed. Preparation and final review follow contribution prompts C and D, with source gaps disclosed. Structural and rendering checks do not establish a new theorem or open status.
\label{end:statement}
\end{document}
